\documentclass[11pt]{article}
\usepackage{mystyle}

\title{Rosen --- \S 9.2: n-ary Relations and Their Applications\\ \large Selected Exercises}
\author{Alston}
\date{Semester 2, 2026}

\begin{document}
\maketitle

% ======================================================================
\begin{exercise}{9}
    The 5-tuples in a 5-ary relation represent these attributes of all
    people in the United States: name, Social Security number, street
    address, city, state.
    \begin{enumerate}[label=(\alph*)]
        \item Determine a primary key for this relation.
        \item Under what conditions would (name, street address) be a
        composite key?
        \item Under what conditions would (name, street address,
        city) be a composite key?
    \end{enumerate}
\end{exercise}

% ======================================================================
\begin{exercise}{18}
    How many components are there in the $n$-tuples in the table
    obtained by applying the join operator $J_3$ to two tables with
    5-tuples and 8-tuples, respectively?
\end{exercise}

% ======================================================================
\begin{exercise}{24}
    Show that if $C$ is a condition that elements of the $n$-ary
    relations $R$ and $S$ may satisfy, then
    \[
    \sigma_C(R - S) = \sigma_C(R) - \sigma_C(S).
    \]
\end{exercise}

% ======================================================================
\begin{exercise}{25}
    Show that if $R$ and $S$ are both $n$-ary relations, then
    \[
    P_{i_1, i_2, \dots, i_m}(R \cup S) = P_{i_1, i_2, \dots, i_m}(R)
    \cup P_{i_1, i_2, \dots, i_m}(S).
    \]
\end{exercise}

% ======================================================================
\begin{exercise}{27}
    Give an example to show that if $R$ and $S$ are both $n$-ary
    relations, then $P_{i_1, i_2, \dots, i_m}(R - S)$ may be different
    from $P_{i_1, i_2, \dots, i_m}(R) - P_{i_1, i_2, \dots, i_m}(S)$.
\end{exercise}

% ======================================================================
\begin{exercise}{32}
    Show that an $n$-ary relation with a primary key can be thought of
    as the graph of a function that maps values of the primary key to
    $(n-1)$-tuples formed from values of the other domains.
\end{exercise}

\end{document}